\section{Đề tài trong một trang}
\label{sec:topic}

\subsection{Bài toán}

Ung thư phổi là nguyên nhân tử vong do ung thư hàng đầu. Nhiều ca lần đầu lộ ra dưới dạng một
\textbf{nốt} hoặc \textbf{khối} trên X-quang ngực. Nốt nhỏ và mờ rất dễ bị bỏ sót.

Mô hình ngôn ngữ--thị giác \en{vision-language model, VLM} là mô hình nhận ảnh và câu hỏi, rồi trả lời bằng chữ.
Người ta đang giới thiệu VLM như ``trợ lý đọc phim''. Nhưng các đánh giá độc lập năm 2026 cho thấy VLM gần như
\textbf{không thấy nốt phổi}. Ví dụ, MedGemma-4B chỉ phát hiện 5,2\% nốt trên bộ JSRT, dù không báo nhầm ca nào
(specificity 100\%). Trong khi đó, phần mềm CNN chuyên dụng đạt AUROC 0,92--0,99. Kiểu sai ``luôn nói bình thường''
này gọi là \textbf{chế độ báo bình thường} \en{normal-mode failure}.

\subsection{Đồ án làm gì}

Đồ án xây một nguyên mẫu trợ lý AI. Trợ lý trả lời: ảnh có nốt/khối không, và nếu có thì ở đâu (bounding box).
Có ba phần việc:

\begin{enumerate}
\item \textbf{Benchmark}: so hai VLM nhỏ (MedGemma-4B, Qwen2.5-VL-3B) với hai CNN (DenseNet-121, YOLOv8) trên cùng cách chia dữ liệu,
      trên ba bộ: VinBigData, NIH ChestX-ray14, JSRT.
\item \textbf{Tinh chỉnh LoRA} MedGemma-4B: chỉ phân loại; phân loại kèm box; có và không cân bằng lớp.
\item \textbf{Kiểm toán ``nhìn đúng chỗ''}: đo IoU, và thử che nốt so với che một vùng ngẫu nhiên cùng diện tích.
\end{enumerate}

\begin{figure}[H]
\centering
\begin{tikzpicture}[node distance=0.5cm and 0.6cm]
\node[box,minimum width=2.6cm] (img) {Ảnh X-quang ngực\\(896$\times$896)};
\node[box,right=1.1cm of img,minimum width=3.0cm] (q) {Câu hỏi cố định:\\\itshape ``Is there a lung nodule\\or mass\ldots? Answer in JSON.''};
\node[gbox,below=0.9cm of $(img)!0.5!(q)$,minimum width=4.5cm] (vlm) {VLM (MedGemma-4B + LoRA)};
\node[file,below=0.8cm of vlm,minimum width=6.4cm] (out) {\{"nodule": "yes", "boxes": [[412, 300, 470, 355]]\}};
\draw[arr] (img) -- (vlm); \draw[arr] (q) -- (vlm); \draw[arr] (vlm) -- (out);
\node[lab,right=0.3cm of out] {đầu ra JSON:\\có/không + box};
\end{tikzpicture}
\caption{Trợ lý nhận ảnh và một câu hỏi cố định, trả lời bằng JSON.}
\label{fig:assistant}
\end{figure}

\subsection{Câu hỏi nghiên cứu và tiêu chí thành công}

\begin{keybox}
\textbf{Q1.} LoRA có đưa sensitivity của VLM, đo ở specificity 0,90, lên \textbf{ít nhất 0,60} trên tập test VinBigData (trong phân phối) không?

\textbf{Q2.} Cải thiện đó có còn khi kiểm tra ngoài phân phối (NIH, JSRT) và khi che nốt không?

\textbf{Tiêu chí thành công} (spec §2): có bảng so sánh 7 cấu hình $\times$ 3 tập test, kèm khoảng tin cậy 95\%,
và kết luận được Q1, Q2 \emph{dù theo hướng nào}. Kết quả âm tính (LoRA không đạt 0,60) vẫn là kết quả hợp lệ.
\end{keybox}

Các khái niệm ``sensitivity ở specificity 0,90'', ``trong/ngoài phân phối'', ``khoảng tin cậy'' được giải thích ở §\ref{sec:background}.

\begin{warnbox}
\textbf{Không viết ``ngang bác sĩ''.} Ngưỡng 0,60 nằm cùng khoảng với sensitivity của người đọc trong vài nghiên cứu,
nhưng các nghiên cứu đó không đo ở specificity 0,90. Spec §2 cấm viết ``ngang bác sĩ''. Đây là \textbf{nguyên mẫu nghiên cứu},
không phải công cụ lâm sàng.
\end{warnbox}
